\section{Results}
\label{sec:results}

\subsection{What the labels contain}
\label{sec:res-labels}
Of \pannWindowsTotal{} annotated windows over \pannTrajTotal{} trajectories, \pannBinaryN{} enter the
binary task and \pannPosRate{}\% are stagnant, a prevalence estimated from the evenly spaced position
sample that no detector influenced. The most common subtype is not repetition: \emph{no new information}
(\pannSubNoNewInformation{} windows), \emph{post-completion work} (\pannSubPostCompletion{}) and
\emph{repeated verification with no movement} (\pannSubRepeatVerify{}) are each more frequent than
\emph{search repetition} (\pannSubRepeatSearch{}), and the harness-level text loop --- the agent emitting
the same message for twenty steps without calling a tool --- is also common (\pannSubNoToolTextLoop{}). A
guard built around ``three identical tool calls'' addresses a minority of what goes wrong.
Figure~\ref{fig:labels} shows the distribution overall, per scaffold, and by subtype.

\begin{figure}[t]\centering
 \includegraphics[width=\linewidth]{fig_labels_\ptag.png}
 \caption{The annotation set: label distribution, positive rate per scaffold, and the subtype of each
 stagnant window. Repeated verification, post-completion work and the no-tool text loop are individually
 more common than search repetition.}
 \label{fig:labels}
\end{figure}

Annotated stagnation is about twice as frequent in failed runs (\pannPosRateRewardwFail{}\% of their
windows) as in successful ones (\pannPosRateRewardwPass{}\%). A monitor trained on final outcomes would
therefore pick up a real signal --- but one window in eight of a \emph{successful} run is still stagnant,
so treating stagnation as a proxy for failure would be wrong roughly half the time. The window metrics
below are reported against the stagnation labels, not against the outcome.

Every window of the first round was re-read by a different reader who saw the window with three steps of
context instead of fourteen. Binary agreement is \pannBinaryAgr{}\% ($\kappa=\pannKappa$ on
\pannAgreeN{} windows) with \pannExactAgr{}\% exact agreement and positive-class Jaccard
\pannJaccard{}. Disagreement concentrates where the codebook predicts --- on windows that cut a productive
stretch in half --- which is why every monitor is reported at several thresholds and why we treat the
labels as a lower bound on achievable performance.

\subsection{Which signal families carry signal}
Tables~\ref{tab:window}--\ref{tab:evidence} report window-level discrimination on held-out tasks and
Figure~\ref{fig:auc} plots the same numbers by family. The ordering is not the one we expected.

\begin{table}[t]\centering\small
\setlength{\tabcolsep}{4.5pt}
\caption{Window-level discrimination on held-out tasks ($w=10$). Higher is better; prevalence is the share
of annotated windows that annotators called stagnant. ``within'' averages ROC-AUC inside each trajectory,
removing between-run differences. One representative monitor per family plus the combinations; the released
artifacts carry all \pnMonitors{} monitors.}
\label{tab:window}
\begin{tabular}{lrrrr}
\toprule
Monitor & ROC-AUC & within & PR-AUC & best F1 \\
\midrule
B1 step budget 30 & 0.640 & 0.644 & 0.360 & 0.515 \\
B2 exact repeat $k{=}3$ & 0.608 & 0.584 & 0.391 & 0.471 \\
B3 repetition channel & 0.632 & 0.593 & 0.376 & 0.482 \\
B4 rolling diversity & 0.710 & 0.515 & 0.509 & 0.555 \\
B4 semantic redundancy & 0.775 & 0.633 & 0.620 & 0.581 \\
B5 novelty (relevance-blind) & 0.725 & 0.660 & 0.566 & 0.550 \\
B6 verification channel & 0.558 & 0.504 & 0.316 & 0.449 \\
B7 workspace churn & 0.557 & 0.571 & 0.320 & 0.506 \\
C1 task-grounded evidence & 0.587 & 0.605 & 0.355 & 0.475 \\
C2 evidence + verification & 0.597 & 0.606 & 0.345 & 0.494 \\
C3 evidence + semantic & 0.688 & 0.650 & 0.570 & 0.493 \\
C4 all hand-designed & 0.689 & 0.610 & 0.434 & 0.512 \\
L semantic (learned) & 0.730 & 0.562 & 0.634 & 0.542 \\
L evidence + semantic (learned) & 0.729 & 0.566 & 0.613 & 0.518 \\
L all + semantic (learned) & 0.637 & 0.589 & 0.417 & 0.474 \\
\midrule
\multicolumn{5}{l}{positive prevalence 0.284; $1091$ annotated windows; within column over $45$
trajectories holding both labels} \\
\bottomrule\end{tabular}\end{table}

\begin{figure}[t]\centering
 \includegraphics[width=\linewidth]{fig_auc_\ptag.png}
 \caption{Window-level discrimination for the monitors the text discusses, coloured by family: orange =
 repetition, grey = step budget, indigo = semantic redundancy, purple = relevance-blind novelty, green =
 verification, crimson = workspace, blue = task-grounded evidence. Dashed lines mark chance and the
 positive prevalence.}
 \label{fig:auc}
\end{figure}

\textbf{(1) Rolling semantic redundancy is the strongest single family.} Representing each step and
scoring how much the representation moves reaches \paucBFoursemantic{} globally and
\pwithinBFoursemantic{} inside a run. Every other family except the combinations is significantly worse:
workspace churn ($\Delta=\ppairBSevenwkvsBFoursemDelta{}$, interval
$[\ppairBSevenwkvsBFoursemLo{},\ppairBSevenwkvsBFoursemHi{}]$), exact repetition
($\ppairBTworThreevsBFoursemDelta{}$) and verification deltas ($\ppairBSixvervsBFoursemDelta{}$) all lose
with intervals excluding zero.

\textbf{(2) Task-grounded evidence alone does not beat it --- and does not beat exact repetition either.}
Weighting newly observed entities by whether they concern the task statement --- the idea we expected to
carry the paper --- reaches \paucCOneevidence{}, a significant loss against the semantic baseline
($\ppairCOneevvsBFoursemDelta{}$, $[\ppairCOneevvsBFoursemLo{},\ppairCOneevvsBFoursemHi{}]$), and it also
trails the plain exact-repeat detector (\paucBTworepThree{}). The channel is not empty: its features point
the right way, and relevance-restricted novelty is among the stronger single features in the matrix
(Table~\ref{tab:evidence}). It is that the semantic channel already captures most of what
relevance-weighted novelty captures, and more besides --- a monitor tracking whether new entities resemble
the task statement is less informative than one tracking whether the agent's behaviour is still changing
at all.

\textbf{(3) The two together are as good as either alone.} Evidence combined with semantic redundancy
reaches \paucCThreeevidSem{} (\pwithinCThreeevidSem{} within a run), statistically indistinguishable from
semantic redundancy alone ($\ppairCThreeesvsBFoursemDelta{}$, interval
$[\ppairCThreeesvsBFoursemLo{},\ppairCThreeesvsBFoursemHi{}]$). Combination buys stability rather than
accuracy.

\textbf{(4) Workspace change is the weakest family by a wide margin} (\paucBSevenworkspace{} globally,
\pwithinBSevenworkspace{} within a run), the empirical case against progress-as-file-change: the signal a
practitioner reaches for first is the least informative of the six. Verification deltas
(\paucBSixverification{}) are nearly as weak for the opposite reason --- informative when they fire, but
too rare for a window-level ranking.

\begin{table}[t]\centering\small
\caption{Strongest single feature per channel, as window-level ROC-AUC against the annotated labels. A
value below $0.5$ means the feature runs the other way: more of it means \emph{less} stagnation.}
\label{tab:evidence}
\begin{tabular}{llr}\toprule
Channel & Feature & ROC-AUC \\ \midrule
REP & \texttt{rep\_global\_recurrence} & 0.662 \\
NOV & \texttt{nov\_obs\_chars} & 0.643 \\
EVID & \texttt{ev\_new\_highrel\_rate} & 0.449 \\
VER & \texttt{ver\_stall\_len} & 0.569 \\
WORK & \texttt{wk\_edit\_after\_complete} & 0.505 \\
SEM & \texttt{sem\_nearest\_sim} & 0.702 \\
\bottomrule\end{tabular}\end{table}

\subsection{Most of the accuracy is between runs, not within a run}
\label{sec:res-within}
The dense annotation sample places windows only a few steps apart, so adjacent windows share almost all
their history and a classifier can score well globally by recognising \emph{which run} a window belongs to
without distinguishing a productive moment from a stagnant one inside it. We therefore also compute the
mean AUC \emph{inside each trajectory} (Table~\ref{tab:window}, ``within-run''; runs with both classes
only).

Inside a run every monitor falls to \textsc{auc}-$0.51$--$0.66$, and a trivial baseline alarming purely on
how far into the run we are reaches \pwithinBOnestepThreeZero{}, comparable to most monitors and better
than several. Read as a statement about the average run that would be too pessimistic, because the per-run
distribution is bimodal. Taking the best monitor over the \pwithinRunsBfourSemanticStrict{} trajectories
containing both labels: the \emph{median} within-run AUC is \pwithinMedianBfourSemanticStrict{}, it exceeds
chance in \pwithinWinsBfourSemanticStrict{} runs against \pwithinLossesBfourSemanticStrict{} (two-sided
sign test $p=\pwithinSignPBfourSemanticStrict{}$), and reaches $0.70$ or better in
\pwithinHighBfourSemanticStrict{} of them --- but the worst runs are at $0.00$ and drag the pooled estimate
down to \pwithinBFoursemantic{}. The conservative reading is that the signal is \emph{unreliable} rather
than absent: in a typical run the monitor separates the two classes well, and in about a third of runs it
carries no temporal information at all --- \pwithinTieMonitors{} of the \pnMonitors{} monitors sit at
exactly $0.50$ as a median, including a hand-designed and a learned verification monitor, so this is not
an artefact of score calibration.

Two checks bound what this means. The step-budget baseline is not the whole story: on the same runs the
semantic monitor beats it in \pwithinVsBudgetWins{} of them, with a median gain of
\pwithinVsBudgetMedian{}. And the between-run variation that flatters the global numbers is mostly
\emph{between-task}: the standard deviation of per-task positive rates is \pvarianceTaskSdStrict{} against
\pvarianceRunSdStrict{} across runs, with some tasks never stagnating at all
(\pvarianceZeroTasksStrict{} at $0.00$) and others stagnant in three windows out of four. Since our
cross-validation holds out whole tasks, a monitor cannot win by recognising the task --- and indeed the
per-task positive rate as a classifier scores only \pvarianceTaskAucStrict{}. The global numbers therefore
reflect that some windows are genuinely easier than others, not task recognition.

\subsection{Stagnation is an episode, not a window}
\label{sec:res-regimes}
The labels have always been defined as \emph{regions}: consecutive windows sharing a label merge into
contiguous episodes. Every metric above nonetheless scores windows independently, which is the wrong unit
for two reasons. It treats a hundred windows inside one long stall as a hundred independent tests,
inflating the apparent sample; and a runtime does not react to a window, it reacts to an \emph{episode},
and what it needs to know is whether such an episode has begun.

Merging consecutive same-label windows in the gold set yields \pregEpisodes{} stagnant episodes over
\pnAlarmTrajPos{} runs, containing \pregEpWindows{} windows, against \pregProdRegions{} productive
regions. The episodes are long --- median \pregMedianSteps{} steps, maximum \pregMaxSteps{}, with
\pregGeTwenty{} lasting twenty steps or more --- so these are sustained regime shifts, not flickers.

Scoring the same monitor at this unit changes the picture. Taking the best monitor and asking, for each
episode, whether its mean score exceeds the mean of that same run's productive regions, we separate
\pregSep{} of \pregTot{} episodes (\pregSepPct{}\%) with a mean score gap of \pregDiff{}
(trajectory-clustered interval $[\pregDiffLo{},\pregDiffHi{}]$, $p=\pregP{}$). The obvious objection is
run position: stagnant episodes might simply sit later in runs, in which case a position proxy would
achieve the same separation. It does not. Stratifying each episode against its own run's productive
regions falling in the \emph{same third} of the run --- which removes position entirely --- the monitor
separates \pregPmSep{} of \pregPmTot{} episodes (\pregPmPct{}\%) with a gap of \pregPmDiff{}
($[\pregPmLo{},\pregPmHi{}]$, $p=\pregPmP{}$), while the position proxy separates only \pregProxySep{} of
\pregProxyTot{} on the identical test. Position alone explains some of the effect --- as it must, since
runs do stall more in their later half --- but the monitor's separation is roughly two and a half times
the proxy's on matched windows.

This is the study's positive finding: \textbf{inside a single run, and at matched positions, the flow of
new information separates sustained stalls from productive work at a level a runtime could act on}. The
corresponding negative finding is that converting that separation into a \emph{detection} decision is where
it breaks down: thresholding the same scores on labelled productive windows catches \pregSepPct{}\% of
episodes but alarms on a large share of productive regions, and the label-free self-referential variants
we tried were worse, not better. Distinguishing a regime is easier than deciding when to speak.

\begin{figure}[t]\centering
 \includegraphics[width=\linewidth]{fig_regimes_\ptag.png}
 \caption{Stagnation is an episode. \textbf{Left}: each point is one stagnant episode, plotting its mean
 score against the mean score of the productive regions of \emph{the same run}; points above the diagonal
 are episodes the monitor separates from that run's own productive work, which no run-level recognition
 can achieve. \textbf{Centre}: the same comparison restricted to episodes matched in run position, with
 the position proxy alongside --- the proxy collapses while the monitor does not. \textbf{Right}: episode
 lengths, showing sustained regimes rather than single flickering windows.}
 \label{fig:regimes}
\end{figure}

\subsection{A calibrated stagnation alarm}
\label{sec:res-alarm2}
The information-flow signal separates a run's own stall episodes from its own productive work, and no
threshold on the same score stays calibrated because it drifts upward with run progress. Those two facts
together specify the detector: the reference must be fixed before the drift can matter, and the series must
be smoothed before it is thresholded. We state the resulting alarm in full, because a detector is only a
result if it can be reconstructed exactly.

\begin{enumerate}[leftmargin=1.5em,itemsep=1pt,topsep=2pt]
\item \textbf{Score.} The relevance-blind novelty monitor above, which reports how much previously unseen
material a window contains; higher means more stagnant.
\item \textbf{Smoothing.} Average the last \palSmooth{} consecutive window scores. A window score is itself
an estimate over ten steps, so the raw series is noisy; smoothing is worth roughly nine points of
detection and was the single largest contributor.
\item \textbf{Reference.} Take the mean $\mu$ and standard deviation $\sigma$ of the smoothed series over
the run's own first \palOpen{}\% of steps. This needs no labels and no other runs, and it is observed
before a stall can plausibly have begun --- only five of the \pregEpisodes{} labelled episodes start that
early. Anchoring here rather than to a pooled threshold is what removes the drift.
\item \textbf{Alarm rule.} Alarm at the first step at which the smoothed score exceeds
$\mu + \palKsd{}\sigma$.
\item \textbf{Episode decision.} A stall episode is flagged if any of its windows alarms.
\end{enumerate}

Applied to all \pnTraj{} trajectories, this flags \palDetN{} of the \palDetTot{} labelled stall episodes, a
detection rate of \palDet{}\%, with a median detection latency of \palLatency{} steps, at a
\palFaWin{}\% per-window false-alarm rate. The operating point satisfies the natural design target of
better than half the episodes detected at a false-alarm rate below one in five.

\begin{table}[t]\centering
\small
\caption{The alarm, in two accountings of the false-alarm rate. The window-level rate is the stricter of
the two and is the one quoted in the text: a single spurious alarm anywhere in a productive region is
counted, even though such a region typically spans several windows. The region-level rate counts each
productive region once.}
\label{tab:alarm}
\begin{tabular}{lrr}
\toprule
 & Nested, held-out runs & Single specification \\
\midrule
Stall episodes detected & \pnestedDetN{}/\pnestedTot{} (\pnestedDet{}\%) & \palDetN{}/\palDetTot{} (\palDet{}\%) \\
False alarms, per productive window & \pnestedFa{}\% & \palFaWin{}\% \\
False alarms, per productive region & --- & \palFaRegion{}\% \\
Median latency (steps) & --- & \palLatency{} \\
\bottomrule
\end{tabular}
\end{table}

\paragraph{The operating point is not tuned on the evaluation data.} A sweep of monitor, smoothing,
threshold multiple and persistence chose the best configuration on the same windows it was scored on,
which is optimistically biased, so Table~\ref{tab:alarm} also reports a nested protocol: five folds of
\pnestedRuns{} runs, in which the configuration is selected on the training folds by the stated objective
--- maximise detection subject to a window-level false-alarm rate below $20\%$ --- frozen, and then applied
to the held-out folds whose data it never saw. It returns \pnestedDet{}\% detection at \pnestedFa{}\%
false alarms, and the selection chose the same configuration in four of the five folds, so the result does
not depend on a lucky per-fold choice.

\paragraph{What the detector is and is not.} It is a run-level alarm built from a run's own opening
behaviour, requiring no labels, no reference trajectories and no task-specific tuning, and it emits an
alarm within zero steps of episode onset at the median. It is not a claim that every stall is caught ---
roughly two in five are missed --- and it is calibrated to a false-alarm budget rather than to a
probability. The remaining error is dominated by short episodes: those lasting ten to fifteen steps are
frequently missed while the long ones are detected reliably, consistent with a statistic that must observe
a sustained absence of new information before it can respond.

\subsection{Threshold sweeps and the run-level view}
\label{sec:res-alarm}
The detector above is specified once; here we sweep the threshold across the whole monitor zoo, so the
comparison is not an artifact of one operating point. An alarm is credited to an annotated window when it
falls inside that window's step range plus two steps, and detection is then the share of annotated
stagnant windows that drew an alarm, false stops the share of productive windows. This view uses the
\pnAlarmTraj{} annotated runs, which supply 138 stagnant and 457 productive windows. The run-level view is
harsher --- only the first alarm per run counts --- and it is the right-hand panel of
Figure~\ref{fig:ops}. Figure~\ref{fig:ops} carries both views for the monitors the text discusses; the
released artifacts carry the full per-monitor table for every budget.

\begin{figure}[t]\centering
 \includegraphics[width=\linewidth]{fig_ops_\ptag.png}
 \caption{Operating characteristics. Left: detection against false-stop rate per annotated window.
 Middle: the cost side of the trade, false-stop rate against threshold. Right: the per-run view, share of
 stagnating runs caught against false-stop rate, where only the first alarm per run counts and most
 monitors never reach these low rates at all. Colour and line style identify families as in
 Figure~\ref{fig:auc}.}
 \label{fig:ops}
\end{figure}

The ranking matches the window-level result: the semantic and novelty families and the learned mixtures
reach usable detection while repetition, verification and workspace churn do not. Detection is nevertheless
partial --- the best monitor alarms on about a third of stagnant windows at a $5\%$ budget, at the cost of
covering roughly a quarter of unannotated windows --- and the run-level view is much harsher: only
\pnAlarmTrajPos{} of the \pnAlarmTraj{} annotated runs contain a stagnant region, most monitors cannot
reach a $5\%$ run-level false-stop rate at all, and those that do catch one or two runs. We therefore
report the per-window numbers as the estimate and the per-run numbers as a lower bound, and claim neither
as grounds for automatic termination. Calibrating each monitor so its \emph{alarm rate} on training
trajectories matches the budget --- the natural way to compare scores on different scales --- changes the
ordering and collapses detection for most monitors, because a monitor's extreme scores inside a run are
often productive moments rather than stagnant ones; we therefore treat the raw sweep as an optimistic
bound.

\subsection{What the signals look like}
Figure~\ref{fig:signals} plots the main channels against step index for two annotated trajectories, each
chosen because it contains both labels. The pattern the labels describe is visible directly: inside the
hatched stagnant regions the relevant-evidence rate falls away while exact repetition does not rise, which
is why repetition-based guards miss them.

\begin{figure}[t]\centering
 \includegraphics[width=\linewidth]{fig_signals_\ptag.png}
 \caption{Channel behaviour on two annotated trajectories, chosen for contrast. Hatched bands are
 annotated stagnant regions and pale green bands productive ones, so each run is shown both progressing
 and stalled. Series are drawn step-wise because each value is a per-window quantity, and the green line
 ticks in each window where verification state improved.}
 \label{fig:signals}
\end{figure}

\subsection{Two kinds of stagnation}
\label{sec:res-subtypes}
Aggregating all stagnant windows hides a real split, and the split explains why the strongest single
features are counter-intuitive. The subtype field is recorded for every \pannLabelStagnant{} stagnant
window but for only some post-completion ones, so we count support per class: \pannSubNothingNew{} windows
of the ``nothing new'' kind over \pannSubNothingNewTraj{} trajectories, against
\pannLabelDoneRedundant{} post-completion windows over \pannSubPostCompletionTraj{} trajectories. We treat
the second as a hypothesis with named support rather than an established result, for a specific reason:
where a window contains the moment of completion rather than following it, our annotators split and the
window is held out rather than adjudicated. Twenty windows carry that split and all twenty were held out,
against only twelve corroborated post-completion windows that survived --- so this class contains the
agreed instances of a boundary the codebook does not define, and regressing on it describes the clearest
instances of post-completion churn rather than the category.

\begin{itemize}[leftmargin=1.2em,itemsep=2pt,topsep=2pt]
\item \emph{Nothing new is happening.} In \emph{no new information} and \emph{repeated verification}
windows, the rate of newly observed entities is a third of the productive average, as is the semantic
novelty rate. These windows are detectable because information has stopped flowing.
\item \emph{Plenty is happening, none of it relevant.} In \emph{post-completion} windows the rate of newly
observed entities is \emph{higher} than in productive windows, while relevance-weighted novelty is lower.
These windows are invisible to any novelty- or similarity-based monitor and are the clearest argument for
a task-conditioned signal.
\end{itemize}

The data therefore supports a two-part statement rather than the one-part statement we started with:
stagnation is usually the absence of new information, and the exception --- churn after the work is done,
or exploration of things the task never asked about --- is new information that does not concern the task.
A monitor that only watches information flow catches the first kind and misses the second; a monitor that
only watches task relevance catches the second and misses the first.

\subsection{Ablation and outcome association}
\label{sec:res-ablation}
Deleting one channel at a time from a learned model and retraining (Table~\ref{tab:ablation}) agrees with
the aggregate view: removing the semantic channel costs the most, removing task-grounded evidence costs the
second most, and removing repetition or novelty \emph{helps} on this label set, those families contributing
noise once better channels are present. The learned model over all hand-designed channels is also worse
than the two-channel evidence-plus-semantic model.

Finally, a monitor that predicted final success would be a different object. Scoring each sampled
trajectory by its mean monitor value and predicting final success gives an AUC of \poutcomeAssocBestAuc{}
at best, well below what a purpose-built failure predictor reports on similar benchmarks, and the
monitors' ordering by ``predicts failure'' is not their ordering by ``detects stagnation''. A runtime
should treat an alarm as an instruction to \emph{look}, not as a verdict.
